%=============================================================================!
% 13.8  WHAT THE TRAJECTORY LOOKS LIKE                                        !
%=============================================================================!
\section{What the trajectory looks like}
\label{sec:th-trajectory}

A thermostat changes what a healthy run looks like: the energy is no
longer constant, and the temperature should fluctuate by a definite
amount. This section shows a healthy thermostatted run of liquid argon
and six faults that thermostats leave, each with the check that catches
it.

\subsection*{A healthy run}

\Cref{fig:th-healthy} follows the liquid under CSVR with $\tau_{\mathrm T}
= \SI{100}{\femto\second}$ for \SI{100}{\pico\second}. From
\SI{5}{\pico\second} on, its kinetic temperature averages
\SI{135.57}{\kelvin}, with a relative spread of 0.0486 against the
canonical 0.0511, and lies within one canonical spread of the target
for 0.719 of the time, close to the 0.683 of a Gaussian
(\cref{sec:sm-probability}). The total energy is no longer kept: it
wanders over \SI{9.17}{\milli\electronvolt} per atom as the thermostat
exchanges heat. The effective energy of \cref{eq:th-effective} stays
within \SI{0.018}{\milli\electronvolt} per atom, five hundred times less,
and falls by \SI{1.6e-5}{\electronvolt\per\pico\second} for the whole
box: it plays the
part the total energy played in \cref{sec:cs-trajectory}. The total
momentum stays below \SI{4e-14}{amu\,\angstrom\per\femto\second}, since
CSVR scales a zero momentum to zero, and the histogram of $K$ follows the
gamma density.

\begin{figure}[htb]
\centering
\includegraphics{ch13_thermostats/healthy}
\caption{A healthy thermostatted run: liquid argon under CSVR with
$\tau_{\mathrm T} = \SI{100}{\femto\second}$ for \SI{100}{\pico\second}.
(a) The kinetic temperature every \SI{100}{\femto\second}, with the
target and one canonical spread either side (dotted). (b) The total
energy and the effective energy per atom, from their starting values. (c)
The kinetic energy over the run, divided by its mean, against the
canonical gamma density (dashed).}
\label{fig:th-healthy}
\end{figure}

\subsection*{Six faults}

\Cref{fig:th-faults} shows six faults and \cref{tab:th-faults} the check
that catches each.

\begin{figure}[p]
\centering
\includegraphics{ch13_thermostats/faults}
\caption{Thermostat faults (teal) against healthy runs (grey dashed). (a)
Liquid argon under Berendsen coupling with $\tau_{\mathrm T} =
\SI{100}{\femto\second}$, against CSVR. (b) The running mean of the
temperature in three runs under Andersen's thermostat at
\SI{0.1}{\per\pico\second}, against three under CSVR. (c) The kinetic
energy of the motion of the centre of mass, started moving, under
Berendsen coupling, against CSVR, on a logarithmic axis. (d) The density
of the potential energy of lithium atoms on the surface, each under its
own Nosé-Hoover thermostat (teal) or Berendsen thermostat (ochre, cut
off), against the exact canonical density. (e) Lithium atoms under CSVR
told $N_{\mathrm f} = 3$: their kinetic temperature over the two freedoms
they have, against CSVR told 2. (f) The mean squared displacement in
liquid argon under Langevin dynamics with $\gamma = \SI{0.1}{\per\pico
\second}$, measured in the cell, against measured from the centre of
mass.}
\label{fig:th-faults}
\end{figure}

\begin{table}[htb]
\centering
\caption{The faults of \cref{fig:th-faults}, the check that catches each,
and its value in the healthy and the broken run; (a) uses the first of
the three runs, so its spreads differ slightly from the three-run means of
\cref{tab:th-compare}.}
\label{tab:th-faults}
\small
\begin{tabular}{@{}llll@{}}
\toprule
fault & check & healthy & broken\\
\midrule
(a) coupling too strong & spread of $T$ / canonical & 0.95 & 0.44\\
(b) coupling too weak & mean $T$ of three runs, K &
   $135.38\pm0.31$ & $131.83\pm1.02$\\
(c) flying ice cube & growth of $\ln K_{\mathrm{com}}$, /ps & $-0.0002$ &
   0.0047\\
(d) one atom, one variable & gap from the exact density & 0.016 & 0.231\\
(e) wrong $N_{\mathrm f}$ & $T$ over the true freedoms, K & 1005 & 1506\\
(f) centre as diffusion & $D$, $10^{-4}$\,\AA$^2$/fs & 3.69 & 4.62\\
\bottomrule
\end{tabular}
\end{table}

\emph{(a) Coupling too strong.} Berendsen coupling with $\tau_{\mathrm T} =
\SI{100}{\femto\second}$ holds the temperature in a band 0.44 as wide as
the canonical one; CSVR with the same time constant gives 0.95. The
average looks perfect, \SI{135.00}{\kelvin}, and only the spread shows
the fault: the check is the relative spread of $T$ against
$\sqrt{2/N_{\mathrm f}}$.

\emph{(b) Coupling too weak.} Andersen's thermostat at
\SI{0.1}{\per\pico\second} lets the temperature wander for tens of
picoseconds: three runs average 133.66, 130.13 and
\SI{131.69}{\kelvin}, counting all $3N$ freedoms, against 135.58, 134.77 and \SI{135.79}{\kelvin}
under CSVR. Nothing in one run looks wrong; the check is the spread of
the run means between independent starts, or the running mean's failure
to settle near the target.

\emph{(c) The flying ice cube.} A velocity of the centre of mass of
\SI{0.0002}{\angstrom\per\femto\second} in each component, left in the
velocities, gives the centre of mass the kinetic energy $K_{\mathrm{com}}$,
whose logarithm grows under Berendsen coupling at
\SI{0.0047}{\per\pico\second}, against the \SI{0.0050}{\per\pico\second}
predicted in \cref{sec:th-berendsen}: $K_{\mathrm{com}}$ rises from
\SI{0.064}{\electronvolt} to \SI{0.161}{\electronvolt} in
\SI{200}{\pico\second}; under CSVR it does not grow. The check is the
kinetic energy of the centre of mass, and of the slowest motions, over a
long run.

\emph{(d) One atom, one variable.} A lithium atom with its own
Nosé-Hoover thermostat misses the exact density of its potential energy
by 0.231 of its peak; with a chain of three, by 0.016; under Berendsen
coupling it is pinned near one energy. The check is a distribution
compared with an exact or independently sampled one, here the density of
$U$; for a large system, the temperatures of its parts
(\cref{sec:sm-trajectory}).

\emph{(e) The wrong number of freedoms in the target.} A thermostat told
three freedoms per lithium atom, which has two, drives each atom's
kinetic energy to $\frac32\kB T_0$, so its true temperature reads
\SI{1506}{\kelvin} for a target of \SI{1000}{\kelvin}. The same fault in
a large system is small but systematic: a thermostat that takes $3N$
freedoms for a run whose centre of mass is held still, as ASE's Berendsen, CSVR and
Nosé-Hoover chain do for argon, raises the temperature of 256 atoms by the
factor $768/765$, \SI{0.5}{\kelvin} at \SI{135}{\kelvin}, about the
scatter between independent runs. For rigid water a thermostat that also
ignores the held distances, as ASE's Nosé-Hoover chain does, raises it by
$576/381 = 1.51$ (\cref{sec:sm-trajectory}); ASE's CSVR and Berendsen
subtract the held distances and miss only the three freedoms of the centre of mass, $384/381$.
The check is the kinetic temperature computed with the freedoms counted
for the run.

\emph{(f) The centre of mass counted as diffusion.} Langevin and Andersen
thermostats do not keep the total momentum, so the centre of mass
wanders; under Langevin dynamics it diffuses, over times long against
$1/\gamma$, with the coefficient $\kB T/(M\gamma)$, which for the whole
box at $\gamma = \SI{0.1}{\per\pico\second}$ is
\SI{1.10e-4}{\angstrom\squared\per\femto\second}
(\cref{ex:th-com}). Displacements measured in the cell include it: the
diffusion coefficient reads $4.62\times10^{-4}$ instead of
$3.69\times10^{-4}$\,\AA$^2$/fs, 1.25 times too large. The check is to measure displacements from the centre of mass.

\begin{practice}
A thermostatted run is read through its kinetic temperature, whose
spread should be $\sqrt{2/N_{\mathrm f}}$ of the target when the freedoms
are counted for the run; through its effective energy, which should
neither wander nor drift; through the kinetic energy of its centre of
mass and of its slowest motions, which should not grow; and, for a
property that depends on the motion, through the same property with the
thermostat weakened, which should not change. Runs from several
independent starts show whether a weakly coupled run has reached its
temperature, a question Chapter~15 answers with error bars.
\end{practice}

\begin{keyidea}
In a healthy thermostatted run the total energy wanders, the effective
energy holds still, the temperature fluctuates by the canonical
$\sqrt{2/N_{\mathrm f}}$, and the centre of mass stays still or is
measured from. Coupling too strong narrows the spread; too weak leaves the
mean unsettled; rescaling and Berendsen coupling fly the ice cube; a single Nosé-Hoover
variable fails on small systems; a miscounted $N_{\mathrm f}$ shifts the
temperature; and a wandering centre of mass inflates diffusion.
\end{keyidea}

\notebook{13}{break the thermostatted liquid in each of the six ways and read the checks}

\begin{selfcheck}
A Langevin run of 500 argon atoms with $\gamma = \SI{1}{\per\pico\second}$
at \SI{135}{\kelvin}: at what rate does its centre of mass diffuse?
\end{selfcheck}
